% ═══════════════════════════════════════════════════════════════════
% Section V: Conclusion
% ═══════════════════════════════════════════════════════════════════

\section{Conclusion}\label{Conclusion}

We introduced CE inference scheduling, a per-slot decision framework that determines whether to execute, approximate, or skip channel estimation given that pilot observations are always available.
The proposed three-tier scheme (Skip/Delta/Full) with an LS-based variation monitor is CE-algorithm-agnostic and requires no modification to the underlying estimator.
Ray-traced evaluations across 5G massive MIMO and 6G ELAA configurations demonstrate up to 98\% skip rate for static users with less than 2\,dB NMSE degradation on unseen base stations.
Computation savings scale with CE cost, confirming that the framework is most valuable in the ELAA regime where advanced estimators dominate the PHY latency budget.

Future work includes integration with GPU-native O-RAN platforms for real-time validation, extension to multi-user scheduling where per-UE skip decisions interact, and combination with temporal channel prediction as an enhanced Delta tier.
